\documentclass[10pt,a4paper]{amsart}
\usepackage{amsmath,amssymb,amsthm,graphicx}
\usepackage{verbatim}
\newenvironment{reference}{\comment}{\endcomment}
%\newenvironment{reference}{}{}
\newenvironment{slogan}{\comment}{\endcomment}
\newenvironment{history}{\comment}{\endcomment}

% For commutative diagrams we use Xy-pic
\usepackage[all]{xy}

% We use 2cell for 2-commutative diagrams.
\xyoption{2cell}
% Theorem environments.
%
\theoremstyle{plain}
\newtheorem{theorem}[subsection]{Theorem}
\newtheorem{proposition}[subsection]{Proposition}
\newtheorem{lemma}[subsection]{Lemma}

\theoremstyle{definition}
\newtheorem{definition}[subsection]{Definition}
\newtheorem{example}[subsection]{Example}
\newtheorem{exercise}[subsection]{Exercise}
\newtheorem{situation}[subsection]{Situation}

\theoremstyle{remark}
\newtheorem{remark}[subsection]{Remark}
\newtheorem{remarks}[subsection]{Remarks}

\numberwithin{equation}{subsection}

% Macros
%
\def\lim{\mathop{\mathrm{lim}}\nolimits}
\def\colim{\mathop{\mathrm{colim}}\nolimits}
\def\Spec{\mathop{\mathrm{Spec}}}
\def\Hom{\mathop{\mathrm{Hom}}\nolimits}
\def\Ext{\mathop{\mathrm{Ext}}\nolimits}
\def\SheafHom{\mathop{\mathcal{H}\!\mathit{om}}\nolimits}
\def\SheafExt{\mathop{\mathcal{E}\!\mathit{xt}}\nolimits}
\def\Sch{\mathit{Sch}}
\def\Mor{\mathop{\mathrm{Mor}}\nolimits}
\def\Ob{\mathop{\mathrm{Ob}}\nolimits}
\def\Sh{\mathop{\mathit{Sh}}\nolimits}
\def\NL{\mathop{N\!L}\nolimits}
\def\CH{\mathop{\mathrm{CH}}\nolimits}
\def\proetale{{pro\text{-}\acute{e}tale}}
\def\etale{{\acute{e}tale}}
\def\QCoh{\mathit{QCoh}}
\def\Ker{\mathop{\mathrm{Ker}}}
\def\Im{\mathop{\mathrm{Im}}}
\def\Coker{\mathop{\mathrm{Coker}}}
\def\Coim{\mathop{\mathrm{Coim}}}

% Boxtimes
%
\DeclareMathSymbol{\boxtimes}{\mathbin}{AMSa}{"02}

%
% Macros for moduli stacks/spaces
%
\def\QCohstack{\mathcal{QC}\!\mathit{oh}}
\def\Cohstack{\mathcal{C}\!\mathit{oh}}
\def\Spacesstack{\mathcal{S}\!\mathit{paces}}
\def\Quotfunctor{\mathrm{Quot}}
\def\Hilbfunctor{\mathrm{Hilb}}
\def\Curvesstack{\mathcal{C}\!\mathit{urves}}
\def\Polarizedstack{\mathcal{P}\!\mathit{olarized}}
\def\Complexesstack{\mathcal{C}\!\mathit{omplexes}}
% \Pic is the operator that assigns to X its picard group, usage \Pic(X)
% \Picardstack_{X/B} denotes the Picard stack of X over B
% \Picardfunctor_{X/B} denotes the Picard functor of X over B
\def\Pic{\mathop{\mathrm{Pic}}\nolimits}
\def\Picardstack{\mathcal{P}\!\mathit{ic}}
\def\Picardfunctor{\mathrm{Pic}}
\def\Deformationcategory{\mathcal{D}\!\mathit{ef}}
\usepackage{cleveref}
\crefname{originallemma}{original lemma}{original lemmas}
\Crefname{originallemma}{Original lemma}{Original lemmas}
\makeatletter
\newlabel{algebra-lemma-completion-flat}{{00MB}{source}{}{sourcepointer-1}{}}
\newlabel{algebra-lemma-completion-flat@cref}{{[originallemma][0][]00MB}{[1][0][]source}{}{}{}}
\makeatother
\crefname{originallemma}{original lemma}{original lemmas}
\Crefname{originallemma}{Original lemma}{Original lemmas}
\makeatletter
\newlabel{algebra-lemma-yet-another-variant-local-criterion-flatness}{{0GEB}{source}{}{sourcepointer-2}{}}
\newlabel{algebra-lemma-yet-another-variant-local-criterion-flatness@cref}{{[originallemma][0][]0GEB}{[1][0][]source}{}{}{}}
\makeatother
\crefname{originalproposition}{original proposition}{original propositions}
\Crefname{originalproposition}{Original proposition}{Original propositions}
\makeatletter
\newlabel{algebra-proposition-Jacobson-permanence}{{00GB}{source}{}{sourcepointer-3}{}}
\newlabel{algebra-proposition-Jacobson-permanence@cref}{{[originalproposition][0][]00GB}{[1][0][]source}{}{}{}}
\makeatother
\crefname{originalreference}{original reference}{original references}
\Crefname{originalreference}{Original reference}{Original references}
\makeatletter
\newlabel{lemma-naively-rig-flat-continuous}{{0GGL}{source}{}{sourcepointer-4}{}}
\newlabel{lemma-naively-rig-flat-continuous@cref}{{[originalreference][0][]0GGL}{[1][0][]source}{}{}{}}
\makeatother
\crefname{originalreference}{original reference}{original references}
\Crefname{originalreference}{Original reference}{Original references}
\makeatletter
\newlabel{lemma-rig-closed-jacobson}{{0GGC}{source}{}{sourcepointer-5}{}}
\newlabel{lemma-rig-closed-jacobson@cref}{{[originalreference][0][]0GGC}{[1][0][]source}{}{}{}}
\makeatother
\crefname{originalreference}{original reference}{original references}
\Crefname{originalreference}{Original reference}{Original references}
\makeatletter
\newlabel{lemma-rig-closed-point-in-localization}{{0GGF}{source}{}{sourcepointer-6}{}}
\newlabel{lemma-rig-closed-point-in-localization@cref}{{[originalreference][0][]0GGF}{[1][0][]source}{}{}{}}
\makeatother
\crefname{originalreference}{original reference}{original references}
\Crefname{originalreference}{Original reference}{Original references}
\makeatletter
\newlabel{lemma-rig-closed-point-relative}{{0GGB}{source}{}{sourcepointer-7}{}}
\newlabel{lemma-rig-closed-point-relative@cref}{{[originalreference][0][]0GGB}{[1][0][]source}{}{}{}}
\makeatother
\title{Naive rig-flatness is rig-flat over a Jacobson base}
\author{The Stacks Project Authors}
\date{stacks-project:restricted.tex:local-snapshot-20261012; GNU FDL 1.2 or later}
\begin{document}
\maketitle
\noindent\textit{Editorial note.} All delivered mathematical text is the author's own, copied from the pinned Stacks Project chapter file restricted.tex; no statement or proof step was rewritten, reordered or removed.\par\medskip
\noindent\textit{Editorial note.} The counted proof is the author's complete argument: the proof is complete, with the localization permanence, the rig-closed lifting and the local flatness criterion all written out\par\medskip
\noindent\textit{Editorial note.} The label lemma-rig-flat-naive names the counted statement, and the proof environment immediately following it in the source is delivered verbatim as its proof.\par\medskip
\noindent\textit{Editorial note.} The counted statement is carried by the author's own lemma environment in the chapter source.\par\medskip
\noindent\textit{Allowed standard prerequisites (editorial).}\begin{itemize}
\item the WAdm\textasciicircum{}Noeth category, localizations A\_\{f\} and their completions A\_\{f\} bracket, and naive rig-flatness on rig-closed primes
\item Jacobson rings and permanence under finite type quotients
\item the rig-closed prime lifting lemmas under localization and contraction
\item the local criterion for flatness with the generating condition on primes
\item finite inclusions of local rings and units
\end{itemize}\medskip
\ifdefined\hypertarget\hypertarget{sourcepointer-1}{}\fi
\noindent\textit{Source reference.} Original lemma 00MB: sources/stacks-chow/tags.txt, source lines 817--817. Prerequisite pointer invoked at proof line(s) 3761; it names a standard fact of a sibling chapter that the counted argument uses, and every step of the counted proof built on it is retained verbatim in the delivered body.\par\medskip
\ifdefined\hypertarget\hypertarget{sourcepointer-2}{}\fi
\noindent\textit{Source reference.} Original lemma 0GEB: sources/stacks-chow/tags.txt, source lines 20471--20471. Prerequisite pointer invoked at proof line(s) 3775; it names a standard fact of a sibling chapter that the counted argument uses, and every step of the counted proof built on it is retained verbatim in the delivered body.\par\medskip
\ifdefined\hypertarget\hypertarget{sourcepointer-3}{}\fi
\noindent\textit{Source reference.} Original proposition 00GB: sources/stacks-chow/tags.txt, source lines 601--601. Prerequisite pointer invoked at proof line(s) 3765; it names a standard fact of a sibling chapter that the counted argument uses, and every step of the counted proof built on it is retained verbatim in the delivered body.\par\medskip
\ifdefined\hypertarget\hypertarget{sourcepointer-4}{}\fi
\noindent\textit{Source reference.} Original reference 0GGL: sources/stacks-restricted/pinned-original-source.tex, source lines 3650--3650. Prerequisite pointer invoked at proof line(s) 3770; it names a standard fact of the same chapter that the counted argument uses, and every step of the counted proof built on it is retained verbatim in the delivered body.\par\medskip
\ifdefined\hypertarget\hypertarget{sourcepointer-5}{}\fi
\noindent\textit{Source reference.} Original reference 0GGC: sources/stacks-restricted/pinned-original-source.tex, source lines 3381--3381. Prerequisite pointer invoked at proof line(s) 3771; it names a standard fact of the same chapter that the counted argument uses, and every step of the counted proof built on it is retained verbatim in the delivered body.\par\medskip
\ifdefined\hypertarget\hypertarget{sourcepointer-6}{}\fi
\noindent\textit{Source reference.} Original reference 0GGF: sources/stacks-restricted/pinned-original-source.tex, source lines 3460--3460. Prerequisite pointer invoked at proof line(s) 3789; it names a standard fact of the same chapter that the counted argument uses, and every step of the counted proof built on it is retained verbatim in the delivered body.\par\medskip
\ifdefined\hypertarget\hypertarget{sourcepointer-7}{}\fi
\noindent\textit{Source reference.} Original reference 0GGB: sources/stacks-restricted/pinned-original-source.tex, source lines 3333--3333. Prerequisite pointer invoked at proof line(s) 3796; it names a standard fact of the same chapter that the counted argument uses, and every step of the counted proof built on it is retained verbatim in the delivered body.\par\medskip
\begin{lemma}
\label{lemma-rig-flat-naive}
Let $\varphi : A \to B$ be an arrow of $\textit{WAdm}^{Noeth}$.
If $A/I$ is Jacobson for some (equivalently any) ideal of definition
$I \subset A$ and $\varphi$ is naively rig-flat, then $\varphi$ is
rig-flat.
\end{lemma}

\begin{proof}
Assume $\varphi$ is naively rig-flat. We first state some obvious
consequences of the assumptions. Namely, let $f \in A$.
Then $A, B, A_{\{f\}}, B_{\{f\}}$
are Noetherian adic topological rings. The maps
$A \to A_{\{f\}} \to B_{\{f\}}$ and $A \to B \to B_{\{f\}}$
are adic and topologically of finite type.
The ring maps $A \to A_{\{f\}}$ and $B \to B_{\{f\}}$
are flat as compositions of $A \to A_f$ and $B \to B_f$
and the completion maps which are flat by
Algebra, Lemma \ref{algebra-lemma-completion-flat}.
The quotients of each of the rings
$A, B, A_{\{f\}}, B_{\{f\}}$ by $I$ is of finite type
over $A/I$ and hence Jacobson too
(Algebra, Proposition \ref{algebra-proposition-Jacobson-permanence}).

\medskip\noindent
Let $\mathfrak q' \subset B_{\{f\}}$ be rig-closed.
It suffices to prove that $(B_{\{f\}})_{\mathfrak q'}$
is flat over $A_{\{f\}}$, see Lemma \ref{lemma-naively-rig-flat-continuous}.
By Lemma \ref{lemma-rig-closed-jacobson} the primes
$\mathfrak q \subset B$ and $\mathfrak p' \subset A_{\{f\}}$
and $\mathfrak p \subset A$ lying under $\mathfrak q'$ are rig-closed.
We are going to apply
Algebra, Lemma \ref{algebra-lemma-yet-another-variant-local-criterion-flatness}
to the diagram
$$
\xymatrix{
B_\mathfrak q \ar[r] &
(B_{\{f\}})_{\mathfrak q'} \\
A_\mathfrak p \ar[u] \ar[r] &
(A_{\{f\}})_{\mathfrak p'} \ar[u]
}
$$
with $M = B_\mathfrak q$.
The only assumption that hasn't been checked yet is the fact
that $\mathfrak p$ generates the maximal ideal of
$(A_{\{f\}})_{\mathfrak p'}$. This follows from
Lemma \ref{lemma-rig-closed-point-in-localization};
here we use that $\mathfrak p$ and $\mathfrak p'$ are rig-closed
to see that $f$ maps to a unit of $A/\mathfrak p$
(this is the only step in the proof that fails without
the Jacobson assumption). Namely,
this tells us that $A/\mathfrak p \to A_{\{f\}}/\mathfrak p'$
is a finite inclusion of local rings
(Lemma \ref{lemma-rig-closed-point-relative})
and $f$ maps to a unit in the second one.
\end{proof}

\end{document}
